\subsection*{Example Code: Using Storage Qualifiers}

The following sketch combines the storage declarations:

\begin{lstlisting}[language=croqtile,caption={Example Code: Using Storage Qualifiers}]
ndims : [20, 15];
local f32 [10, 10] d0;
shared f16 [ndims] d1;        // Shared storage for intermediate data across threads
global f32 [100, 200] matrix; // Global storage for large dataset
// Perform computations on these arrays, utilizing the appropriate storage space
dma.copy d0 => d1;            // Copy from local to shared storage
dma.copy d1 => matrix;        // Copy from shared to global storage
\end{lstlisting}

The declarations illustrate three storage locations. The copy lines sketch the direction of movement between them. To turn this sketch into a working transfer, select source and destination views with compatible shapes: the full arrays declared here have different sizes. A storage specifier chooses a location; it does not make differently shaped buffers interchangeable.
D1 is an intermediate array placed in shared storage to be used by multiple threads within a block. Matrix is a large dataset placed in global storage, where it can be accessed by all threads.
